La \textbf{planificación del proyecto} se ha estructurado teniendo en cuenta las fases principales del \textbf{ciclo de vida del software} y la naturaleza académica del \textit{Trabajo Fin de Grado}. Al tratarse de un proyecto realizado por un \textbf{único desarrollador}, la planificación combina actividades de análisis, diseño, construcción, pruebas y documentación, distribuidas a lo largo del periodo de realización del trabajo.

\subsection{Calendario general}

El calendario se presenta por meses, ya que el proyecto se desarrolla como un \textit{Trabajo Fin de Grado} individual y las actividades no se ejecutan como tareas aisladas de corta duración, sino como fases de \textbf{análisis}, \textbf{diseño}, \textbf{construcción}, \textbf{validación} y \textbf{documentación} que se solapan parcialmente. Esta granularidad permite reflejar la evolución real del trabajo sin introducir una falsa precisión diaria que no se corresponde con la \textbf{naturaleza incremental} del proyecto.

El proyecto se ha organizado en \textbf{fases sucesivas}, aunque con cierto solapamiento entre actividades. En particular, la documentación se ha ido actualizando de forma progresiva conforme avanzaban el análisis, el diseño y la construcción del sistema. La Tabla~\ref{tab:project-plan-schedule} resume esta \textbf{planificación general}.

\begin{table}[htbp]
    \centering
    \small
    \caption{Planificación general del proyecto}
    \label{tab:project-plan-schedule}
    \begin{tabularx}{\textwidth}{|L{3.5cm}|L{2.8cm}|Y|}
        \hline
        \textbf{Fase} & \textbf{Periodo aproximado} & \textbf{Actividades principales} \\
        \hline
        \textbf{Inicio y definición del proyecto} & Febrero de 2026 & Definición inicial del problema, contacto con el tutor, delimitación del alcance y preparación de la toma de requisitos. \\
        \hline
        \textbf{Captura de requisitos} & Marzo de 2026 & Reunión con el cliente, análisis de la situación actual, identificación de procesos, problemas y necesidades funcionales. \\
        \hline
        \textbf{Análisis del sistema} & Marzo - abril de 2026 & Definición de objetivos, requisitos funcionales, no funcionales, requisitos de información, reglas de negocio y división modular. \\
        \hline
        \textbf{Diseño del sistema} & Abril de 2026 & Diseño de arquitectura, diseño conceptual y físico de datos, diseño de interfaz y organización técnica de la solución. \\
        \hline
        \textbf{Construcción} & Abril - mayo de 2026 & Preparación del entorno, implementación de la base funcional, autenticación, roles, usuarios y estructura inicial de la aplicación. \\
        \hline
        \textbf{Pruebas y validación} & Mayo de 2026 & Validación manual de funcionalidades, comprobaciones de ejecución, revisión de flujos principales y pruebas en distintos dispositivos. \\
        \hline
        \textbf{Documentación y cierre} & Mayo - junio de 2026 & Redacción de la memoria, anexos, manual del desarrollador, manual de usuario, revisión final y preparación de la entrega. \\
        \hline
    \end{tabularx}
\end{table}
\subsection{Diagrama de Gantt}

La Figura~\ref{fig:project-plan-gantt} muestra una representación temporal resumida de las principales fases del proyecto. El \textit{Diagrama de Gantt} refleja la \textbf{naturaleza incremental} del trabajo, en la que algunas tareas, especialmente la \textbf{documentación}, se desarrollan de forma paralela a las fases técnicas.

\begin{figure}[htbp]
    \centering
    \footnotesize

    \begin{ganttchart}[
            hgrid,
            vgrid,
            x unit=0.85cm,
            y unit title=0.65cm,
            y unit chart=0.72cm,
            title height=1,
            bar height=0.45,
            milestone height=0.45,
            title label font=\footnotesize\bfseries,
            bar label font=\footnotesize,
            milestone label font=\footnotesize,
            bar/.append style={
                draw=black,
                fill=white
            },
            milestone/.append style={
                draw=black,
                fill=black
            },
            bar label node/.append style={
                text width=4.2cm,
                align=right,
                anchor=east,
                xshift=-0.15cm
            },
            milestone label node/.append style={
                text width=4.2cm,
                align=right,
                anchor=east,
                xshift=-0.15cm
            }
        ]{1}{10}

        \gantttitle{2026}{10} \\
        \gantttitle{Feb}{2}
        \gantttitle{Mar}{2}
        \gantttitle{Abr}{2}
        \gantttitle{May}{2}
        \gantttitle{Jun}{2} \\

        \ganttbar{Inicio y definición}{1}{2} \\
        \ganttbar{Captura de requisitos}{3}{4} \\
        \ganttbar{Análisis del sistema}{3}{6} \\
        \ganttbar{Diseño del sistema}{5}{6} \\
        \ganttbar{Construcción}{5}{8} \\
        \ganttbar{Pruebas y validación}{7}{8} \\
        \ganttbar{Documentación y cierre}{7}{10} \\
        \ganttmilestone{Entrega final}{9}

    \end{ganttchart}

    \caption{Diagrama de Gantt resumido del proyecto}
    \label{fig:project-plan-gantt}
\end{figure}

\subsection{Hitos principales}

Los hitos principales considerados durante el desarrollo han sido los siguientes:

\begin{itemize}
    \item \textbf{Aceptación y definición inicial} del \textit{Trabajo Fin de Grado}.
    \item \textbf{Realización de la reunión de toma de requisitos} con el cliente.
    \item \textbf{Definición de la división modular} del sistema.
    \item \textbf{Elección del marco tecnológico principal}.
    \item \textbf{Diseño de la arquitectura y del modelo de datos}.
    \item \textbf{Implementación de la base funcional} de la aplicación.
    \item \textbf{Validación de los flujos principales} desarrollados.
    \item \textbf{Finalización de la memoria} y preparación de los entregables.
\end{itemize}

\subsection{Comparación entre esfuerzo planificado y esfuerzo real}

La Tabla~\ref{tab:project-plan-effort} recoge una \textbf{comparación aproximada} entre el esfuerzo inicialmente estimado y el esfuerzo finalmente dedicado a las principales actividades del proyecto. Dado que el desarrollo se ha realizado de forma individual y no se ha utilizado un sistema formal de registro horario, los valores reales se han estimado a partir de la evolución del proyecto, los entregables generados y el tiempo dedicado a cada fase.

\begin{table}[htbp]
    \centering
    \small
    \setlength{\tabcolsep}{4pt}
    \caption{Comparación entre esfuerzo planificado y esfuerzo real}
    \label{tab:project-plan-effort}
    \begin{tabularx}{\textwidth}{|L{5.7cm}|>{\centering\arraybackslash}X|>{\centering\arraybackslash}X|>{\centering\arraybackslash}X|}
        \hline
        \textbf{Actividad} & \textbf{\footnotesize Planificado (h)} & \textbf{\footnotesize Real aprox. (h)} & \textbf{\footnotesize Desviación (h)} \\
        \hline
        Análisis de requisitos & 30 & 32 & +2 \\
        \hline
        Diseño de arquitectura & 25 & 27 & +2 \\
        \hline
        Diseño de base de datos & 20 & 20 & 0 \\
        \hline
        Diseño de interfaz & 30 & 28 & -2 \\
        \hline
        Desarrollo \textit{backend} & 90 & 95 & +5 \\
        \hline
        Desarrollo \textit{frontend} & 100 & 102 & +2 \\
        \hline
        Integración de servicios externos & 20 & 12 & -8 \\
        \hline
        Pruebas y validación & 40 & 36 & -4 \\
        \hline
        Documentación del proyecto & 35 & 48 & +13 \\
        \hline
        \textbf{Total} & \textbf{390} & \textbf{400} & \textbf{+10} \\
        \hline
    \end{tabularx}
\end{table}

La \textbf{desviación global} ha sido moderada. El incremento de esfuerzo se concentra principalmente en la \textbf{documentación} y en la \textbf{consolidación de la base funcional del sistema}. Por otro lado, algunas integraciones externas previstas se han mantenido como \textit{líneas de trabajo futuro}, reduciendo el esfuerzo real dedicado a esa fase dentro del alcance final del \textit{Trabajo Fin de Grado}.
